\documentclass{article}
% Source: Topic_01_Teacher_Edition.pdf, PDF pages 12-23 (printed pages 5A-12B)
\usepackage{amsmath}
\usepackage{amssymb}
\usepackage{graphicx}
\usepackage{tikz}
\usepackage{pgfplots}
\pgfplotsset{compat=1.18}
\usepackage{booktabs}
\usepackage{xcolor}

\newenvironment{lesson-meta}{}{}        % lesson code, title, objective, EQ, vocab
\newenvironment{item-analysis}{}{}      % Savvas's Example × DOK table
\newenvironment{model-discuss}{}{}      % Launch opener
\newenvironment{concept-box}{}{}        % in-lesson concept reference
\newenvironment{concept-summary}{}{}    % end-of-lesson summary
\newenvironment{example}[1]{}{}         % #1 = example number
\newenvironment{tryit}[1]{}{}           % #1 = tryit number
\newenvironment{practice}[2]{}{}        % #1 = practice number, #2 = DOK
\newenvironment{te-addendum}[2]{}{}     % #1 = anchor example, #2 = type
\newcommand{\answer}[1]{}               % answer key, inside any env
\newcommand{\placeholder}[2]{}          % #1 = type, #2 = description
\newcommand{\uncertain}[2]{#1}          % #1 = best guess, #2 = alternatives
\newcommand{\te}[1]{}                   % teacher-only inline annotation

\begin{document}


% Source page 12: 5A Lesson Overview
\begin{lesson-meta}
lesson: 1-1
title: Key Features of Functions
standards: none printed
practices: none printed
objective: I can interpret key features of linear, quadratic, and absolute value functions given an equation or a graph.

essential question: How do key features help you understand functions and their graphs?

\textbf{VOCABULARY}
\item average rate of change
\item interval notation
\item maximum
\item minimum
\item set-builder notation
\item zero of a function
\textbf{TOPIC VOCABULARY}
\te{No topic vocabulary list is printed on these lesson pages. The I CAN statement is on PDF page 13; the Essential Question is on PDF page 14. No numbered content-standard or Mathematical Practice codes are printed on the overview or student-page header.}
\begin{tabular}{ll}
Lesson Code & 1-1 \\
Title & Key Features of Functions \\
Standards & none printed \\
I CAN Objective & I can interpret key features of linear, quadratic, and absolute value functions given an equation or a graph. \\
Essential Question & How do key features help you understand functions and their graphs? \\
Lesson Vocabulary & average rate of change; interval notation; maximum; minimum; set-builder notation; zero of a function \\
Topic Vocabulary & none printed \\
\end{tabular}
\end{lesson-meta}

\begin{te-addendum}{0}{GeneralizeETP}
\textbf{Lesson Overview: FOCUS}
\textbf{Objective}
Students will be able to:
\item Identify key features of a graph of a function, including the intercepts, positive and negative intervals, and areas where the function is increasing or decreasing.
\item Calculate and interpret the average rate of change of a function over a specified interval.
\item Write the domain and range of functions using set-builder and interval notations.
\textbf{Essential Understanding}
The key features of a graph—including the domain, range, and intercepts—reveal the relationship between two quantities.
\textbf{COHERENCE}
Previously in earlier courses, students:
\item Created, graphed, and identified key features of functions.
In this lesson, students:
\item Identify and interpret the key features of the graph of a function, including the domain, range, intercepts, and areas where the graph is increasing or decreasing.
Later in this topic, students will:
\item Use the key features of the graph of a function to identify transformations and write equations.
\textbf{RIGOR}
This lesson emphasizes a blend of conceptual understanding and application.
\item Students identify the intercepts of a function and interpret them in the context of real-world problems, such as finding the distance that a car can travel before it runs out of gas.
\end{te-addendum}

\begin{te-addendum}{0}{ELLAddendum}
\textbf{Vocabulary Builder}
Glossary is available at SavvasRealize.com.
\textbf{REVIEW VOCABULARY}
\begin{tabular}{ll}
English & Spanish \\
average rate of change & tasa promedio de cambio \\
interval notation & notación de intervalo \\
maximum & máximo \\
minimum & mínimo \\
set-builder notation & notación de conjuntista \\
zero of the function & cero de una función \\
\end{tabular}
\textbf{VOCABULARY ACTIVITY}
Prepare students for understanding interval notation and set-builder notation by reviewing the terms domain, range, maximum, and minimum. Explain that set-builder notation uses a verbal description or an inequality to describe domain and range and that interval notation is a set of real numbers represented by the pair of numbers that are its minimum and maximum boundaries. Have students identify each of the following as set-builder notation or interval notation.
(\{x\mid x\text{ is a real number}\}) blank. \answer{set-builder notation}
(-\infty,5) blank. \answer{interval notation}
\te{Printed the symbols “−∞, 5” without visible enclosing delimiters; likely an interval whose endpoint delimiters are missing in the source. The bare parentheses above follow house math style and do not resolve the missing source delimiters.}
([-2,\infty)) blank. \answer{interval notation}
(\{y\mid y<1\}) blank. \answer{set-builder notation}
\end{te-addendum}

\begin{te-addendum}{0}{LearnTogether}
\textbf{Student Companion}
Students can do their work for the lesson in their Student Companion or on Savvas Realize.
\placeholder{illustration}{Student Companion cover: dark background, multicolored loops, white text “Student Companion” at top and “enVision Algebra 2” near bottom.}
\end{te-addendum}

\begin{te-addendum}{0}{GeneralizeETP}
\textbf{Mathematics Overview}
Students find and interpret key features of linear, quadratic, absolute value, and constant functions, including the x- and y-intercepts, the positive and negative intervals, and intervals where the function is increasing and decreasing.
Students use the x-intercepts of a function to determine the zeros of a function, and they use set-builder and interval notation to state domain and range.
\textbf{Applying Math Practices}
\textbf{Construct Viable Arguments}
Students use stated mathematical assumptions about the minimum value of a function to construct an argument explaining why the function has no maximum value.
\textbf{Attend to Precision}
Students demonstrate understanding of the meanings of mathematical symbols when they use information about the amount of gas in a car’s tank and the rate at which it is being used to describe the key features of the graph that models the situation.
\te{SavvasRealize.com: Program Resources.}
\end{te-addendum}

% Source page 13: 5B STEP 1 Explore
\begin{te-addendum}{0}{LearnTogether}
\textbf{EXPLORE \& REASON}
\textbf{INSTRUCTIONAL FOCUS} Students explore the relationship between the depth of the diver and the time since she started the dive to understand how one affects the other. This leads to understanding the key features of functions.
\textbf{STUDENT COMPANION} Students can complete the Explore \& Reason activity in their Student Companion.
\end{te-addendum}
\begin{te-addendum}{0}{GeneralizeETP}
\textbf{Before—WHOLE CLASS}
\textbf{Implement Tasks that Promote Reasoning and Problem Solving ETP}
Q: What do you notice about the graph?
\answer{The graph decreases from left to right until the point (5,-8), and then the graph increases.}
\textbf{During—SMALL GROUP}
\textbf{Support Productive Struggle in Learning Mathematics ETP}
Q: What do the coordinates of a point on the graph tell you about the position of the diver?
\answer{The y-coordinate tells you the depth of the diver, and the x-coordinate tells you how many seconds have passed since the start of the dive.}
Q: What does a positive value of (y) tell you about the diver?
\answer{A positive value of (y) tells you that the diver is above the surface of the water.}
\textbf{For Early Finishers}
Q: Draw the graph of (y=|x-11|-10) to represent another diver. What is the starting position of the diver? What is the deepest point of the dive? What is the average speed of the diver?
\answer{The starting position is (0,1). The deepest point of the dive is (-10), or 10 feet below the surface. The average speed of the diver in the water is 1 ft/s.}
\textbf{After—WHOLE CLASS}
\textbf{Facilitate Meaningful Mathematical Discourse ETP}
Q: Are there x-values of the graph that do not make sense in the context of the problem?
\answer{Yes; values less than 0 do not make sense because (x=0) represents the time the dive starts.}
Q: Are there y-values of the graph that do not make sense in the context of this problem?
\answer{Yes; y-values greater than the starting point of 2 do not make sense. The diver will not continue the path once she breaks the surface of the water.}
\end{te-addendum}
\begin{te-addendum}{0}{HabitsOfMind}
\textbf{Use with EXPLORE \& REASON}
Q: Reason What do the points where the graph intersects the x-axis tell you about the dive?
\answer{The points on the x-axis are when the diver is at the surface.}
\end{te-addendum}
\begin{model-discuss}
\textbf{EXPLORE \& REASON}
A diver is going through ocean search-and-rescue training. The graph shows the relationship between her depth and the time in seconds since starting her dive.
\placeholder{graph}{Diver graph in ocean illustration. Horizontal axis t, “Time (s),” numbered 1 through 9 at unit spacing, with origin 0; vertical axis d, “Depth (ft),” labels 2, 0, -2, -4, -6, -8 at two-foot spacing. Visible window approximately t=0 to 9 and d=-8 to 2, with arrowheads on positive t and both vertical ends. Gray dashed V-shaped path from (0,2) through (1,0) to a filled black point (5,-8), then up through (9,0). The vertex callout reads “5 s, −8 ft.” Boat at upper left and illustrated diver on ascending right segment; sea surface coincides with d=0. No other endpoint circle is visibly marked.}
A. What details can you determine about the dive from the coordinates of the point (5,-8)?
\answer{The diver reached the lowest point of the dive, 8 ft below the surface of the water, 5 s after beginning the dive.}
B. What is the average speed of the diver in the water? How can you tell from the graph?
\answer{2 ft/s; Sample: There is a slope of 2 ft/s from (5,-8) to (9,0).}
C. Which point on the graph shows the starting location of the diver? Explain.
\answer{(0,2); This is the point when the time is 0 s and the dive begins.}
D. Communicate Precisely What does the V-shape of the graph tell you about the dive? What information does it not tell you about the dive?
\answer{The diver dove down and then swam back up. It does not tell you if the diver moved horizontally or how far she moved.}
\te{The page prints the activity twice, once as a digital screenshot and once under STUDENT EDITION. Both contain the same prompt and graph. The digital screenshot includes “Go back,” “Enter your answer.” fields, and the EXPLORE \& REASON heading. The response text above is the printed SAMPLE STUDENT WORK. Explore \& Reason is available at SavvasRealize.com.}
\end{model-discuss}

% Source page 14: 5 STEP 2 Understand & Apply
\begin{te-addendum}{0}{LearnTogether}
\textbf{INTRODUCE THE ESSENTIAL QUESTION}
\textbf{Establish Mathematics Goals to Focus Learning ETP}
Introduce students to the key features of linear, quadratic, and absolute value functions. Explain that interpreting the key features of the graphs of these functions can help solve problems about functions.
\te{Examples are available at SavvasRealize.com. The top of the student-page image repeats part D of the Explore \& Reason activity from PDF page 13.}
\end{te-addendum}
\begin{concept-box}
\textbf{ESSENTIAL QUESTION} How do key features help you understand functions and their graphs?
\textbf{CONCEPT Set-Builder Notation and Interval Notation}
Set-builder notation is used to describe the elements of a set. A variable represents any number in the set, followed by a description of the elements of the set.
(\{y\mid y\text{ is an integer and }y\geq10\})
\te{Callout: The set of all numbers (y) such that (y) is an integer and is greater than or equal to 10.}
Interval notation is used to describe a continuous set of real numbers, or interval, by the minimum and maximum values.
((-5,0]) is the same as (\{y\mid-5<y\leq0\}).
\te{Callout: Use parentheses when the minimum or maximum is not included in the set. Use square brackets when it is included.}
([-3,\infty)) is the same as (\{y\mid-3\leq y\}).
((-\infty,7)) is the same as (\{y\mid y<7\}).
\te{Callout: Use (\pm\infty) if there is no minimum or no maximum value in the set.}
\end{concept-box}
\begin{te-addendum}{1}{PurposefulQuestions}
\textbf{Additional Examples—ADDITIONAL EXAMPLE 1}
Q: What are the domain and range of the function (y=-x^2+3)?
\textbf{Domain:}
You can square any real number, so any number can be input for (x). The domain is all real numbers, or ((-\infty,\infty)).
\textbf{Range:}
The maximum value of the function is (y=3). So, the range is all real numbers less than or equal to 3, or ((-\infty,3]).
\answer{Domain: all real numbers, or ((-\infty,\infty)); range: ((-\infty,3]).}
\placeholder{graph}{Orange downward parabola y=-x^2+3, vertex (0,3), y-intercept 3; horizontal x and vertical y axes with arrowheads. Square grid approximately -5 to 5 on each axis, unit grid spacing; visible x labels -4, -2, 0, 4 and y labels 4, 2, -2, -4. Curve has arrows on both descending ends. No individual points or x-intercepts are labeled.}
\textbf{Example 1} Help students find the domain and range of a quadratic function with a negative leading coefficient.
Q: Can values in the range be as large as you want them to be?
\answer{No. The function has a maximum, which is equal to 3. Values in the range of (f) cannot be greater than 3.}
\end{te-addendum}
\begin{te-addendum}{3}{PurposefulQuestions}
\textbf{ADDITIONAL EXAMPLE 3}
Q: For what intervals is (y=-x^2+9) positive? For what intervals is the function negative?
Drag the slider to graph the function.
\answer{The function is positive: ((-3,3)). The function is negative: ((-\infty,-3)), ((3,\infty)). The function is neither positive nor negative at the x-intercepts of (-3) and 3.}
\placeholder{graph}{Desmos screenshot: x and y axes with arrowheads, x labels -10, -5, 0, 5, 10, y labels -5, 5, 10, 15; unit grid lines and approximately x=-12 to 12, y=-9 to 16. A downward parabola has a blue portion above the x-axis and orange portions below it. The displayed vertex is approximately (0,5), and a lower-left slider labeled c displays 5. No individual points are labeled.}
\te{Printed equation and answer describe (y=-x^2+9), but the screenshot slider reads 5 and the displayed vertex is approximately (0,5); likely the screenshot shows a different slider state.}
\textbf{Example 3} Help students find the intervals over which a quadratic function with a negative leading coefficient is positive and negative.
Q: How do the positive and negative intervals of (f(x)=-x^2+4) compare to the positive and negative intervals of (g(x)=x^2-4)?
\answer{They are opposite. The positive intervals of (f) are negative intervals of (g), and vice versa.}
\te{Additional Example 3 is Powered by Desmos. Activities are available at SavvasRealize.com. Both Additional Example screenshots include “Go back,” “ESSENTIAL QUESTION,” and “SOLUTION.”}
\end{te-addendum}

% Source page 15: 6 STEP 2 Understand & Apply
\begin{te-addendum}{1}{PurposefulQuestions}
\textbf{EXAMPLE 1 Understand Domain and Range}
\textbf{Pose Purposeful Questions ETP}
\textbf{PART A}
Q: What difference do you notice between set-builder notation and interval notation?
\answer{Set-builder notation uses a verbal description or an inequality with the variable listed. Interval notation uses a pair of values that have a left minimum boundary and a right maximum boundary, with no variable listed.}
\textbf{PART B}
Q: How would the domain and range be different if the function were given without real-life context?
\answer{If the function (f(x)=400x) were given without the context, the domain and range would be all real numbers.}
\end{te-addendum}
\begin{example}{1}
\textbf{Understand Domain and Range}
A. What are the domain and range of the function defined by (y=x^2-3)?
The set of all possible inputs for a relation is called the domain.
The set of all possible outputs for a relation is called the range.
\placeholder{graph}{Red upward parabola y=x^2-3, vertex (0,-3). Square grid with x and y axes, unit spacing and window approximately [-5,5] by [-5,5]. Both axes have arrows. Labels on each axis are -4, -2, 0, 2, 4. The curve has upward arrows at both upper ends. No points or x-intercepts are labeled.}
\te{Callout: This graph represents a function because each input has exactly one output.}
\te{Callout: In (y=x^2-3), (x) can be any real number, while (y) cannot be less than 3.}
\te{Printed “cannot be less than 3”; likely “cannot be less than −3,” consistent with the printed equation, graph, and range.}
Using set-builder notation, the domain of this function is (\{x\mid x\text{ is a real number}\}) and in interval notation the domain is ((-\infty,\infty)).
Using set-builder notation, the range of the function is (\{y\mid y\geq-3\}) and in interval notation the range is ([-3,\infty)).
\answer{A. Domain: (\{x\mid x\text{ is a real number}\}), ((-\infty,\infty)); range: (\{y\mid y\geq-3\}), ([-3,\infty)).}
\te{STUDY TIP: An interval with excluded boundary points is called “open” and is represented by open circle end points on a number line. An interval with included boundary points is called “closed” and is represented by solid circle end points on a number line.}
B. An airtanker flies over forest fires and drops water at a constant rate until its tank is empty. What are the domain and range of the function that represents the volume of water the airtanker can drop in (x) seconds?
\placeholder{photo}{Yellow airtanker dropping water above a forest fire. Callouts: “Rate: 400 gals. per second” and “8,000 gals. in the tank.”}
The function is (f(x)=400x). The airtanker cannot drop water for a negative number of seconds, so (x\geq0). The tanker can drop water for a maximum of (\frac{8,000}{400}=20\text{ s}) before running out of water, so (x\leq20).
The domain is (\{x\mid0\leq x\leq20\}), or ([0,20]).
The airtanker cannot drop a negative number of gallons, and its maximum capacity is 8,000 gal.
The range is (\{y\mid0\leq y\leq8,000\}), or ([0,8,000]).
\answer{B. Domain: (\{x\mid0\leq x\leq20\}), or ([0,20]); range: (\{y\mid0\leq y\leq8,000\}), or ([0,8,000]).}
\end{example}
\begin{tryit}{1}
What are the domain and range of each function? Write the domain and range in set-builder notation and in interval notation.
a. (y=|x-4|)
b. (y=6x-2x^2)
\answer{a. domain: (\{x\mid x\text{ is a real number}\}); ((-\infty,\infty)); range: (\{y\mid y\geq0\}); ([0,\infty)). b. domain: (\{x\mid x\text{ is a real number}\}); ((-\infty,\infty)); range: (\{y\mid y\leq4.5\}); ((-\infty,4.5]).}
\end{tryit}
\begin{te-addendum}{1}{ElicitEvidence}
\textbf{Elicit and Use Evidence of Student Thinking ETP}
Q: How can a graph help you find the domain and range of a function?
\answer{When you graph the function, you can see the general behavior and shape of the graph, including its minimum and maximum values. Using what you see, you can translate the information into set-builder or interval notation.}
\te{Examples are available at SavvasRealize.com.}
\end{te-addendum}


% Source page 16: 7 STEP 2 Understand & Apply
\begin{te-addendum}{2}{PurposefulQuestions}
\textbf{EXAMPLE 2 Find x- and y-Intercepts}
\textbf{Pose Purposeful Questions ETP}
Q: Why would you want to find the x- and y-intercepts of a function?
\answer{Knowing the x- and y-intercepts can help you graph the function.}
Q: What do you know about the coordinates of the x- and y-intercepts?
\answer{The x-coordinate of a y-intercept is 0, and the y-coordinate of an x-intercept is 0.}
\end{te-addendum}
\begin{te-addendum}{2}{LearnTogether}
\textbf{Learn Together: Offer Help to Others}
Doing math and thinking about math can be challenging. Ask: What can you say to help others feel confident? What can you do to help others persevere?
\end{te-addendum}
\begin{example}{2}
\textbf{Find x- and y-intercepts}
\textbf{APPLICATION}
A. A car starts a journey with a full tank of gas. The equation (y=16-0.05x) relates the number of gallons of gas, (y), left in the tank to the number of miles the car has traveled, (x). What are the x- and y-intercepts of the graph of this equation, and what do they represent about the situation?
\placeholder{graph}{Gas remaining versus travel distance: horizontal x-axis “Travel Distance (mi),” labels 0, 80, 160, 240, 320, 360, grid spacing 40 miles; vertical y-axis “Gas Remaining (gal),” labels 4, 8, 12, 16, 20 and grid spacing 2 gallons. Both axes have arrowheads. Gray descending line segment y=16-0.05x from (0,16) to (320,0), with no marked endpoint circles. Adjacent callouts define both intercepts.}
\placeholder{photo}{Rear view of a gray car on a curved two-lane road in a rocky desert landscape; photo underlies the intercept callouts next to the graph.}
\te{Callout: A y-intercept is the y-coordinate of a point where a graph intersects the y-axis.}
\te{Callout: An x-intercept is the x-coordinate of a point where a graph intersects the x-axis.}
\te{LEARN TOGETHER: Do you seek help when needed? Do you offer help and support others?}
The graph above intersects the x-axis at (320,0), so the x-intercept is 320. This means that the car can travel 320 mi before it runs out of gas.
The graph intersects the y-axis at (0,16), so the y-intercept is 16. This means the car has 16 gal of gas when it starts its trip.
\answer{A. x-intercept: 320; the car can travel 320 mi before it runs out of gas. y-intercept: 16; the car has 16 gal of gas when it starts its trip.}
B. What are the x- and y-intercepts of the graph of (y=|x|-3)?
Find the x-intercept(s) algebraically:
(y=|x|-3)
(0=|x|-3)
(3=|x|)
(x=\pm3)
\te{Callout: The y-coordinate of an x-intercept is 0.}
\te{Callout: Both 3 and (-3) are 3 units away from 0, so they both have an absolute value of 3.}
The x-intercepts of the graph of (y=|x|-3) are (-3) and 3. The x-intercepts are also the zeros of the function because they are the input values that result in a function output value of 0.
Find the y-intercept algebraically:
(y=|x|-3)
(y=|0|-3)
(y=0-3)
(y=-3)
\te{Callout: The x-coordinate at the y-intercept is 0.}
The y-intercept of the graph of (y=|x|-3) is (-3).
\answer{B. x-intercepts: (-3) and 3; y-intercept: (-3).}
\end{example}
\begin{tryit}{2}
What are the x- and y-intercepts of (g(x)=4-x^2)?
\answer{x-intercepts: (-2) and 2; y-intercept: 4}
\end{tryit}
\begin{te-addendum}{2}{ElicitEvidence}
\textbf{Elicit and Use Evidence of Student Thinking ETP}
Q: What do you know about the x-intercepts of the function by looking at the function?
\answer{The function is a quadratic function, so it may have 2 zeros.}
\end{te-addendum}
\begin{te-addendum}{2}{CommonError}
\textbf{Common Error}
\textbf{Try It! 2} Some students may find only one x-intercept, 2. Have students check their work by substituting 0 for (y) and solving for (x): (x^2=4) and (x=\pm\sqrt{4}), or 2 and (-2). So there are two x-intercepts, 2 and (-2).
\end{te-addendum}
\begin{te-addendum}{2}{HabitsOfMind}
\textbf{Use with EXAMPLES 1 \& 2}
Q: Make Sense and Persevere A function does not have any x-intercepts. What might be true about its domain and range?
\answer{The domain could be any set of real numbers, but the range cannot include 0.}
\end{te-addendum}
\begin{te-addendum}{2}{RtISupport}
\textbf{Support Struggling Students}
\textbf{USE WITH EXAMPLE 2A} Students may have trouble identifying the x- and y-intercepts in a real-world context.
\item Have students practice finding the x- and y-intercepts of linear functions.
Q: What are the x- and y-intercepts of (g(x)=2x-1)?
\answer{x-intercept: (\frac{1}{2}); y-intercept: (-1)}
Q: What are the x- and y-intercepts of (h(x)=-3x+6)?
\answer{x-intercept: 2; y-intercept: 6}
Q: What are the x- and y-intercepts of (j(x)=4x)?
\answer{x-intercept: 0; y-intercept: 0}
\te{Examples are available at SavvasRealize.com.}
\end{te-addendum}

% Source page 17: 8 STEP 2 Understand & Apply
\begin{te-addendum}{3}{PurposefulQuestions}
\textbf{EXAMPLE 3 Identify Positive or Negative Intervals}
\textbf{Pose Purposeful Questions ETP}
Q: What do you notice about the graph of the function where the positive and negative intervals meet?
\answer{The boundary points of the intervals are the x-intercepts of the function.}
\end{te-addendum}
\begin{example}{3}
\textbf{Identify Positive or Negative Intervals}
For what intervals of (x) is (f(x)=x^2-9) positive? For what intervals of (x) is the function negative?
Use technology to graph the function:
\placeholder{graph}{Upward parabola f(x)=x^2-9 with vertex (0,-9), x-intercepts -3 and 3. Curve is blue where above the x-axis and red where below, labeled f beside the lower right red branch. Unit x grid, y grid spacing 2; x labels -4, -2, 0, 2, 4, y labels 8, 4, -4, -8. Window approximately x=-6 to 6, y=-12 to 12. Axes and both upward curve ends have arrows. No point circles.}
\te{Callout: (f(x)>0) (above the x-axis) when (x<-3) and when (x>3).}
\te{Callout: (f(x)<0) (below the x-axis) between the x-intercepts of (-3) and 3.}
The function is positive at ((-\infty,-3)) and ((3,\infty)).
The function is negative at ((-3,3)).
The function is neither positive nor negative at the x-intercepts of (-3) and 3.
\answer{Positive: ((-\infty,-3)) and ((3,\infty)); negative: ((-3,3)); neither positive nor negative at the x-intercepts of (-3) and 3.}
\te{COMMON ERROR: Be careful not to confuse a positive function value and a positive rate of change. A positive rate of change means the y-values of the function are increasing but are not necessarily greater than 0.}
\end{example}
\begin{tryit}{3}
a. For what interval(s) is (h(x)=2x+10) positive?
b. For what interval(s) is (h(x)=2x+10) negative?
\answer{a. positive: ((-5,\infty)). b. negative: ((-\infty,-5)).}
\end{tryit}
\begin{te-addendum}{3}{ElicitEvidence}
\textbf{Elicit and Use Evidence of Student Thinking ETP}
Q: How can you determine the intervals over which a function with a connected graph is positive or negative without graphing the function?
\answer{Find the x-intercepts of the function. Choose an x-value between consecutive x-intercepts and evaluate the function at that value. If the value is positive, the function is positive on the interval between the x-intercepts you selected. If the value is negative, the function is negative on that interval.}
\end{te-addendum}
\begin{te-addendum}{4}{PurposefulQuestions}
\textbf{EXAMPLE 4 Identify Where a Function Increases or Decreases}
\textbf{Pose Purposeful Questions ETP}
Q: Can there be a function that only increases or only decreases?
\answer{Yes; a linear function does not have a maximum or minimum, so the function will not change from increasing to decreasing.}
Q: What do you know about the values of a function on either side of a maximum?
\answer{The values to the left of the maximum are increasing, and the values to the right are decreasing.}
\end{te-addendum}
\begin{example}{4}
\textbf{Identify Where a Function Increases or Decreases}
For what values of (x) is (g(x)=2-|x|) increasing? For what values is it decreasing?
Construct a table and sketch a graph to represent the function.
\begin{tabular}{ll}
(x) & (g(x)) \\
(-3) & (-1) \\
(-2) & 0 \\
(-1) & 1 \\
0 & 2 \\
1 & 1 \\
2 & 0 \\
3 & (-1) \\
\end{tabular}
\placeholder{graph}{Red inverted V g(x)=2-|x| with vertex (0,2), intercepts (-2,0) and (2,0), label g along right branch. Cartesian axes x and y with arrows, unit square grid; x labels -6, -4, -2, 0, 2, 4, 6, y labels -4, -2, 4; window approximately x=-7 to 7 and y=-6 to 6. Red arrows at both descending ends. Blue directional arrows above the branches point up and right on the left branch and down and right on the right branch. No point circles.}
\te{Callout: The greatest value a function attains is the maximum of the function. The least value a function attains is the minimum.}
\te{Callout: (g(x)) is decreasing from 0 to (\infty).}
\te{Callout: (g(x)) is increasing from (-\infty) to 0.}
The values of (g(x)) are increasing on the x-interval ((-\infty,0)).
The values of (g(x)) are decreasing on the x-interval ((0,\infty)).
\answer{Increasing: ((-\infty,0)); decreasing: ((0,\infty)).}
\te{REASON: Intervals of the domain explain what is happening in the range of the function. Increasing or decreasing is determined as the x-values increase or as you trace the function from left to right. Why do you think the boundaries are open?}
\end{example}
\begin{tryit}{4}
For what values of (x) is each function increasing? For what values of (x) is it decreasing?
a. (f(x)=x^2-4x)
b. (f(x)=-2x-3)
\answer{a. increasing: ((2,\infty)); decreasing: ((-\infty,2)). b. not increasing on any interval; decreasing: ((-\infty,\infty)).}
\end{tryit}
\begin{te-addendum}{4}{HabitsOfMind}
\textbf{Use with EXAMPLES 3 \& 4}
Q: Use Structure Can a function be increasing and negative on the same interval? Explain.
\answer{Yes; the graph of the function can go from lower left to upper right, but still lie below the x-axis.}
\end{te-addendum}
\begin{te-addendum}{4}{RtIExtend}
\textbf{Extend Student Thinking}
\textbf{USE WITH EXAMPLE 4} Have students compare two functions on an interval and describe the functions as increasing or decreasing.
Q: Compare (f(x)=x^2+4) and (g(x)=-|x|+1) on the interval ((-\infty,0)). Describe the functions as increasing or decreasing.
\answer{(f(x)=x^2+4) is decreasing on this interval. (g(x)=-|x|+1) is increasing on this interval.}
\te{Examples are available at SavvasRealize.com.}
\end{te-addendum}

% Source page 18: 9 STEP 2 Understand & Apply
\begin{te-addendum}{5}{GeneralizeETP}
\textbf{EXAMPLE 5 Understand Average Rate of Change Over an Interval}
\textbf{Build Procedural Fluency From Conceptual Understanding ETP}
Q: What do you know about the average rate of change over any interval of a linear function?
\answer{The average rate of change over any interval is equal to the slope of the function.}
Q: What does it mean for a function to have different average rates of change over different intervals?
\answer{A function with different average rates of change over different intervals has varying rates of increase (or decrease). For example, between 4 and 10 minutes, the cost of a ride increases, on average, by \$0.50 each minute. That increase is less than rides lasting between 10 and 18 minutes, where the cost increases by \$1.50 each minute, on average.}
\end{te-addendum}
\begin{example}{5}
\textbf{Understand Average Rate of Change Over an Interval}
\textbf{CONCEPTUAL UNDERSTANDING}
A. What is the average rate of change of a function (f(x)) over the interval ([a,b])?
The interval starts at a value of (f(a)) when (x=a) and ends at a value of (f(b)) when (x=b).
The total change in the function values is (f(b)-f(a)).
The length of the interval is (b-a).
The average rate of change is the ratio (\frac{f(b)-f(a)}{b-a}). This is the same as the slope of the line segment between the points ((a,f(a))) and ((b,f(b))).
\placeholder{graph}{Schematic increasing red curved graph labeled f(x) on x/y axes with arrowheads. No numeric scale or ticks. Two filled blue points labeled (a,f(a)) and (b,f(b)); a and b marked on x-axis, both positive. Purple dashed secant segment joins the points. Green dashed horizontal run labeled b-a and blue dashed vertical rise labeled f(b)-f(a) form a right triangle. The red curve continues with arrows to the left and upward to the right.}
\answer{A. (\frac{f(b)-f(a)}{b-a}), the slope of the line segment between ((a,f(a))) and ((b,f(b))).}
B. A rideshare program charges fees based on how many minutes a ride lasts. What do the average rates of change over the time intervals ([0,4]), ([10,20]), and ([20,24]) indicate about the cost of rides?
\placeholder{graph}{Rideshare cost graph: x-axis “Time (min)” from 0 to 28, numbered 4, 8, 12, 16, 20, 24, 28; y-axis “Cost (dollars)” from 0 to 28, numbered 4, 8, 12, 16, 20, 24, 28. Grid spacing 2 on both axes. Filled green points labeled (4,6), (10,9), (18,21), (20,25), (24,26). Red horizontal segment from (0,6) to (4,6), gray straight segment to (10,9), purple increasing curved segment through (18,21) to (20,25), blue straight segment through (24,26) extending to the right boundary near (28,27). Axes have arrowheads; curve ends have no visible arrowheads.}
On ([0,4]), the average rate of change is (\frac{6-6}{4-0}=0) meaning the cost does not change on the interval. Any ride 4 minutes or less will cost \$6.
On ([10,20]), the average rate of change is (\frac{25-9}{20-10}=\frac{16}{10}) which is equivalent to the ratio \$1.60:1 minute. For each additional minute of the ride on this interval, the cost will increase an average of \$1.60.
On ([20,24]), the average rate of change is (\frac{26-25}{24-20}=\frac{1}{4}); at this stage of the ride, the cost increases an average of 25 cents per extra minute.
So the company charges the greatest rate for rides lasting between 10 and 20 minutes, and it gives riders a significant price break for taking rides lasting more than 20 minutes.
\answer{B. On ([0,4]), no change in cost: any ride 4 minutes or less costs \$6. On ([10,20]), the cost increases an average of \$1.60 per additional minute. On ([20,24]), the cost increases an average of 25 cents per extra minute.}
\te{LOOK FOR RELATIONSHIPS: If ((x_1,y_1)) and ((x_2,y_2)) are points on the graph of the linear function (y=mx+b), then the average rate of change in the interval ([x_1,x_2]) is (m=\frac{y_2-y_1}{x_2-x_1}).}
\end{example}
\begin{tryit}{5}
What do the average rates of change of the function (y=|x|+2) over the intervals ([-2,0]), ([0,3]), and ([-2,3]) indicate about the function?
\answer{rates of change: (-1); 1; (\frac{1}{5}); Since the rates of change vary, the function is nonlinear.}
\end{tryit}
\begin{te-addendum}{5}{HabitsOfMind}
\textbf{Use with EXAMPLE 5}
Q: Construct Arguments If a function has a positive average rate of change over an interval, does that mean that the function must be increasing over that interval? Explain.
\answer{No; a function can change from increasing to decreasing inside the interval but still have a positive average rate of change over the interval if the final function value is greater than the initial function value.}
\end{te-addendum}
\begin{te-addendum}{5}{ELLAddendum}
\textbf{English Language Learners (Use with EXAMPLE 5)}
\textbf{SPEAKING—BEGINNING} In pairs, have students discuss the similarities and differences between the terms line, segment, and line segment.
Q: What is a line?
\answer{a straight set of points that extends forever in opposite directions}
Q: What is a segment?
\answer{a part of a line with two endpoints and all points between them}
Q: Do the terms segment of a line and line segment have the same meaning?
\answer{The terms mean the same thing.}
\textbf{WRITING—DEVELOPING} In their journals, have students write the definition of the prefix inter- (between). Then have them write the answers to the following questions.
Q: What are some examples of intervals of time?
\answer{hours, days, weeks, months}
Q: What is an interval in music?
\answer{the distance between two musical notes}
Q: How are intervals used in math?
\answer{They define a range of values over which a function is studied.}
Q: What other words do you know that start with the prefix inter-?
\answer{interlude, intercept, interfere, intervene, interim}
\textbf{READING—EXPANDING} In pairs, have students take turns reading the example title, the problem statement, and the paragraph where average rate of change is defined. As they read, have them point to the terms average rate of change, ratio, and slope.
Q: What is the ratio described?
\answer{the change of y-values to the change of x-values}
Q: How is the average rate of change similar to the slope of a line?
\answer{It is the same ratio between points.}
Q: How is the average rate of change different from the slope of a line?
\answer{The slope of a line is constant over any interval of the line. The average rate of change for a nonlinear graph may change, based on the interval.}
\te{Examples are available at SavvasRealize.com.}
\end{te-addendum}


% Source page 19: 10 STEP 2 Understand & Apply
\begin{te-addendum}{ConceptSummary}{PurposefulQuestions}
\textbf{CONCEPT SUMMARY Some Functions With Key Features}
Q: What do these functions have in common?
\answer{The domain of each of these functions is all real numbers. Each function has a y-intercept.}
\end{te-addendum}
\begin{te-addendum}{ConceptSummary}{CommonError}
\textbf{Do You UNDERSTAND? | Do You KNOW HOW?}
\textbf{Common Error}
\textbf{Exercise 4} Students may forget when to use a bracket and when to use a parenthesis in interval notation. Have students make flashcards to review. One side should have interval notation; the other side should represent the same interval on a number line. Students should focus on recognizing that a parenthesis in the interval notation corresponds to an open circle on the number line (or positive/negative infinity), and a bracket represents a closed circle.
\end{te-addendum}
\begin{concept-summary}
\textbf{CONCEPT SUMMARY Some Functions With Key Features}
\begin{tabular}{lllll}
FUNCTION & Linear (y=x) & Quadratic (y=x^2) & Absolute Value (y=|x|) & Constant (y=1) \\
KEY FEATURES: Domain & ((-\infty,\infty)) & ((-\infty,\infty)) & ((-\infty,\infty)) & ((-\infty,\infty)) \\
Range & ((-\infty,\infty)) & ([0,\infty)) & ([0,\infty)) & (\{y\mid y=1\}) \\
Increasing & ((-\infty,\infty)) & ((0,\infty)) & ((0,\infty)) & \\
Decreasing & & ((-\infty,0)) & ((-\infty,0)) & \\
Positive & ((0,\infty)) & (\{x\mid x\neq0\}) & (\{x\mid x\neq0\}) & ((-\infty,\infty)) \\
Negative & ((-\infty,0)) & never & never & never \\
Intercepts & The x-intercept is 0. The y-intercept is 0. & The x-intercept is 0. The y-intercept is 0. & The x-intercept is 0. The y-intercept is 0. & There is no x-intercept. The y-intercept is 1. \\
\end{tabular}
\placeholder{graph}{Summary Linear graph: red line y=x from lower left to upper right, arrowheads both ends. Square unit grid approximately [-4,4] by [-4,4], x and y axes with arrows, origin O and labels -2, 2 on both axes. No point circles.}
\placeholder{graph}{Summary Quadratic graph: red upward parabola y=x^2 with vertex at (0,0), both upper branches arrowed. Square unit grid approximately [-4,4] by [-4,4], x/y axes with arrows, origin O, -2 and 2 labels on both axes. Black callout line from “vertex” points to origin.}
\placeholder{graph}{Summary Absolute Value graph: red V y=|x|, vertex at (0,0), arrowheads on upper left and upper right ends. Square unit grid approximately [-4,4] by [-4,4], x/y axes with arrows, origin O, -2 and 2 labels on both axes. Black callout line from “vertex” points to origin.}
\placeholder{graph}{Summary Constant graph: red horizontal line y=1 with arrows at both ends. Square unit grid approximately [-4,4] by [-4,4], x/y axes with arrows, origin O, -2 and 2 labels on both axes. No points or callouts.}
\textbf{Do You UNDERSTAND?}
1. ESSENTIAL QUESTION How do key features help you understand functions and their graphs?
\answer{Key features give you an overall description for the graph. The basic shape is determined by the type of function. Domain and range can give you the orientation as well as finding the x and y intercepts. Knowing intervals where functions are increasing, decreasing, positive and negative also help you graph the function.}
2. Vocabulary Define the term zero of a function in your own words.
\answer{The zeros of a function are the input values that result in a function output value of 0.}
3. Error Analysis Lonzell said the function shown in the graph is positive on the interval ((-1,5)) and negative on the interval ((-5,-1)). Identify and correct Lonzell’s error.
\placeholder{graph}{Exercise 3: red V-shaped function with vertex (-1,-3), x-intercepts -4 and 2, and y-intercept -2. Square unit grid approximately x=-6 to 6 and y=-4 to 4. Axes x/y with arrows, origin O, x labels -4, -2, 2, 4 and y labels -2, 2, 4. Red arrowheads at upper left and upper right curve ends. No point circles.}
\answer{Lonzell confused positive with increasing and negative with decreasing. The function is positive on the intervals ((-\infty,-4)) and ((2,\infty)) and negative on the interval ((-4,2)).}
\textbf{Do You KNOW HOW?}
Find each key feature.
\placeholder{graph}{Exercises 4–11 share a red polyline with open circle (-4,1), vertex (-2,-1), vertex (2,3), and filled endpoint (4,0), joined by straight segments. Intercepts at (-3,0), (-1,0), (4,0), and (0,1). Unit square grid approximately x=-5 to 4 and y=-2 to 4; x/y axes arrowed, origin O; x labels -4, 2, 4 and y label 2. No curve arrowheads.}
4. domain
\answer{((-4,4])}
5. range
\answer{([-1,3])}
6. x-intercept(s)
\answer{(x=-3), (x=-1), (x=4)}
7. y-intercept(s)
\answer{(y=1)}
8. interval(s) where the graph is positive
\answer{((-4,-3)) and ((-1,4))}
9. interval(s) where the graph is decreasing
\answer{((-4,-2)) and ((2,4])}
10. interval(s) where the graph is increasing
\answer{((-2,2))}
11. average rate of change on ([-1,4])
\answer{0}
\te{Concept Summary, Do You Understand, and Do You Know How are available at SavvasRealize.com.}
\end{concept-summary}


% Source page 20: 11 STEP 3 Practice & Problem Solving
\begin{te-addendum}{Practice}{GeneralizeETP}
\textbf{PRACTICE \& PROBLEM SOLVING}
\textbf{Lesson Practice} You may opt to have students complete the automatically scored Practice and Problem Solving items online powered by MathXL for School.
Choose from: Lesson Practice; Adaptive Practice.
You may also take advantage of the bank of exercises for assigning additional practice.
\textbf{Assignment Guide}
\begin{tabular}{ll}
On-Level & Advanced \\
12, 13, 15, 18--33 & 12--33 \\
\end{tabular}
\textbf{Engage Through Student Choice}
Promote student agency by allowing students to choose practice items. You may structure this choice in many ways.
For example:
Assign each section a point value. Students choose at least one item from each section and items chosen should have a minimum of 20 total points.
\begin{tabular}{ll}
Understand, Apply & 2 points each \\
Practice & 1 point each \\
Assessment Practice & 1 point each \\
Performance Task & 3 points \\
\end{tabular}
\te{Practice \& Problem Solving and video tutorials are available at SavvasRealize.com. Scan the QR code to access a video tutorial.}
\end{te-addendum}
\begin{item-analysis}
\begin{tabular}{lll}
\toprule
Example & Items & DOK \\
\midrule
1 & 18, 23 & 1 \\
2 & 19, 24, 32 & 1 \\
2 & 30 & 3 \\
3 & 20, 25 & 1 \\
3 & 12 & 2 \\
3 & 13 & 3 \\
4 & 21, 26, 31 & 1 \\
4 & 14--17, 28 & 3 \\
5 & 22, 27 & 1 \\
5 & 29, 33 & 3 \\
\bottomrule
\end{tabular}
\end{item-analysis}
\begin{practice}{12}{2}
UNDERSTAND. Reason The graph of (y=-\frac{1}{2}x+2) is negative over the interval ((4,\infty)) and positive over the interval ((-\infty,4)). What happens on the graph when (x=4)? Explain.
\answer{There is a zero at (x=4) because the function is changing from negative to positive at (x=4).}
\te{Printed “changing from negative to positive”; likely “changing from positive to negative” as x increases, consistent with the printed equation and prompt.}
\end{practice}
\begin{practice}{13}{3}
Error Analysis Describe and correct the error a student made in finding the interval(s) over which the function is positive and negative.
\placeholder{graph}{Student-work error graphic on pale blue square grid: red downward parabola with vertex approximately (1,4) and x-intercepts -1 and 3; y-intercept approximately 3. Gray horizontal and vertical axes with arrowheads but no printed axis names, numbers, or tick labels. Red arrows on both descending curve ends. Student labels below: “positive: [-1,3]” and “negative: (-∞,-1) and (3,∞)”; large red X marks the work incorrect.}
\answer{The zeros, or values along the x-axis, are neither positive nor negative. These values should not be included in the interval over which the function is positive. The function is positive on the interval ((-1,3)) and negative on the intervals ((-\infty,-1)) and ((3,\infty)).}
\end{practice}
\begin{practice}{14}{3}
Use Structure Sketch a graph given the following key features.
\begin{tabular}{ll}
domain: ((-4,4)) & range: ((-4,6]) \\
increasing: ((-4,1)) & decreasing: ((1,4)) \\
x-intercepts: (-2,0), (3,0) & y-intercept: (0,4) \\
negative: ((-4,-2)) and ((3,4)) & positive: ((-2,3)) \\
\end{tabular}
\answer{See graph.}
\placeholder{graph}{Printed answer on PDF page 21: red polyline from open endpoint (-4,-4) through (-2,0) and (0,4) to peak (1,6), then through (3,0) to open endpoint (4,-3). No arrows on polyline. Unit square grid, visible window approximately x=-5 to 5 and y=-4 to 6. Axes x/y arrowed, origin O; x labels -4, -2, 2, 4 and y labels -2, 2, 4.}
\te{The printed answer for this item is on PDF page 21.}
\end{practice}
\begin{practice}{15}{3}
Construct Arguments A student says that all linear functions are either increasing or decreasing. Do you agree? Explain.
\answer{No; linear functions can also be constant, which means neither increasing nor decreasing.}
\te{Answer printed on PDF page 21.}
\end{practice}
\begin{practice}{16}{3}
Higher Order Thinking A relative maximum of a function occurs at the highest point on a graph over a certain interval. A relative minimum of a function occurs at the lowest point on a graph over a certain interval. Explain how to identify a relative maximum and a relative minimum of a function using key features.
\answer{A relative maximum occurs when a function changes from increasing to decreasing. A relative minimum occurs when a function changes from decreasing to increasing.}
\te{Answer printed on PDF page 21.}
\end{practice}
\begin{practice}{17}{3}
Model With Mathematics For a graph of speed in miles per hour as a function of time in hours, what does it mean when the function is increasing? Decreasing?
\answer{When the graph is increasing, the speed is increasing. When the graph is decreasing, the speed is decreasing.}
\end{practice}
\begin{practice}{18}{1}
PRACTICE. Use the graph of the function for Exercises 18--22.
\placeholder{graph}{Shared graph for Exercises 18–22: red upward parabola with vertex (-1,-9), x-intercepts -4 and 2, y-intercept -8. Unit grid, visible window approximately x=-5 to 5 and y=-10 to 2; x-axis labels -3, -1, O, 1, 3; y-axis labels -2, -4, -6, -8. Axes x/y have arrowheads; both upper curve ends have red arrows. No plotted-point circles.}
Identify the domain and range of the function. SEE EXAMPLE 1
\answer{domain: ((-\infty,\infty)); range: ([-9,\infty))}
\end{practice}
\begin{practice}{19}{1}
Identify the x- and y-intercepts of the function. SEE EXAMPLE 2
\te{Use the shared graph printed above Exercise 18.}
\answer{x-intercepts: (-4), 2; y-intercept: (-8)}
\end{practice}
\begin{practice}{20}{1}
On what intervals is the function positive? On what intervals is it negative? SEE EXAMPLE 3
\te{Use the shared graph printed above Exercise 18.}
\answer{negative: ((-4,2)); positive: ((-\infty,-4)) and ((2,\infty))}
\end{practice}
\begin{practice}{21}{1}
On what intervals is the function increasing? On what intervals is it decreasing? SEE EXAMPLE 4
\te{Use the shared graph printed above Exercise 18.}
\answer{decreasing: ((-\infty,-1)); increasing: ((-1,\infty))}
\end{practice}
\begin{practice}{22}{1}
What is the average rate of change over the interval ((-3,2))? SEE EXAMPLE 5
\te{Use the shared graph printed above Exercise 18.}
\answer{1}
\end{practice}
\begin{practice}{23}{1}
Use the graph of the function for Exercises 23--27.
\placeholder{graph}{Shared graph for Exercises 23–27: blue polyline joins filled endpoint (-5,0) to vertex (-1,2) to filled endpoint (5,-1). The descending segment crosses y-axis at 1.5 and x-axis at 3. No line arrows. Unit square grid approximately [-5,5] by [-5,5]; axes x/y arrowed, origin O; x labels -4, -2, 2 and y labels -4, -2, 2, 4.}
Identify the domain and range of the function. SEE EXAMPLE 1
\answer{domain: ([-5,5]); range: ([-1,2])}
\end{practice}
\begin{practice}{24}{1}
Identify the x- and y-intercepts of the function. SEE EXAMPLE 2
\te{Use the shared graph printed above Exercise 23.}
\answer{x-intercepts: (-5), 3; y-intercept: 1.5}
\end{practice}
\begin{practice}{25}{1}
Determine over what interval the function is positive or negative. SEE EXAMPLE 3
\te{Use the shared graph printed above Exercise 23.}
\answer{positive: ((-5,3)); negative: ((3,5])}
\end{practice}
\begin{practice}{26}{1}
Determine over what interval the function is increasing or decreasing. SEE EXAMPLE 4
\te{Use the shared graph printed above Exercise 23.}
\answer{increasing: ([-5,-1)); decreasing: ((-1,5])}
\end{practice}
\begin{practice}{27}{1}
What is the average rate of change over the interval ((-1,5))? SEE EXAMPLE 5
\te{Use the shared graph printed above Exercise 23.}
\answer{(-\frac{1}{2})}
\end{practice}

% Source page 21: 12 STEP 3 Practice & Problem Solving
\begin{te-addendum}{Practice}{GeneralizeETP}
\textbf{Answers}
\te{The answers to Exercises 14–16 printed at the top left of this page are attached to their corresponding practice blocks above. Practice \& Problem Solving is available at SavvasRealize.com.}
\end{te-addendum}
\begin{practice}{28}{3}
APPLY. Communicate Precisely Martina is filling an empty (100\text{ ft}^3) container with sand at a rate of (1.25\text{ ft}^3/\text{min}). Describe and interpret the key features of the graph of the amount of sand inside the container.
\answer{The origin is the x and y intercept and informs you the container is empty before it starts filling with sand. The domain is ([0,80]) which are the times when the container is filling with sand. The range is ([0,100]) which are the amounts of sand measured in cubic feet inside the container while it is filling.}
\end{practice}
\begin{practice}{29}{3}
Make Sense and Persevere The graph shows a jumper’s height, (y), in feet (x) seconds after getting onto a trampoline.
\placeholder{graph}{Green repeated downward-curved bounce arcs over a trampoline illustration with two jumpers. Horizontal x-axis “Time (Seconds),” labels 0.5, 1, 1.5, 2, 2.5 with small grid spacing approximately 0.125 second; vertical y-axis “Height (ft),” labels -0.5, 0.5, 1, 1.5 with small grid spacing approximately 0.125 foot. Origin O; axes arrowed. Curve starts at (0,-0.5), rises to a maximum near (0.375,1.75), falls to (0.75,-0.5), rises to (1.125,1.75), falls to (1.5,-0.5), rises to (1.875,1.75), falls to (2.25,-0.5). No individually marked points, endpoint circles, or curve arrows. Perspective illustration distorts the grid slightly.}
a. What are the x- and y-intercepts? Explain what the x- and y-intercepts represent.
\answer{a. The x-intercepts are approximately 0.05, 0.7, 0.8, 1.45, 1.55, and 2.2. The x-intercepts represent the moment the jumper is at the height of the frame of the trampoline. The y-intercept is (-0.5) and represents the sag of the trampoline when the jumper pushes down at the start of the first jump.}
b. Over what intervals is the graph positive? Explain what the positive intervals represent.
\answer{b. ((0.05,0.7)), ((0.8,1.45)), and ((1.55,2.2)); The positive intervals represent the times when the jumper is above the frame of the trampoline.}
c. Over what intervals is the graph negative? Explain what the negative intervals represent.
\answer{c. ((0,0.05)), ((0.7,0.8)), ((1.45,1.55)), and ((2.2,2.25)); The negative intervals represent the times when the jumper is below the frame of the trampoline.}
d. Choose an interval. Find the average rate of change on the interval and interpret its meaning.
\answer{d. Sample answer ([0.75,1.125]); 6 feet per second; The average rate of change over the interval ([0.75,1.125]) represents the average speed of the jumper on the upward part of a jump.}
\end{practice}
\begin{practice}{30}{3}
Model With Mathematics Yanaye starts playing a game on her cell phone with the battery fully charged, and plays until the phone battery dies. While playing the game, the charge in Yanaye’s battery decreases by half a percent per minute.
a. Write a function for the percent charge in the battery while Yanaye is playing the game.
\answer{a. (f(x)=100-0.5x)}
b. What is the domain and range of the function?
\answer{b. domain: ([0,200]); range: ([0,100])}
c. How long can Yanaye play the game?
\answer{c. 3 hours 20 minutes}
\end{practice}
\begin{practice}{31}{1}
ASSESSMENT PRACTICE. Given the graph, select yes or no for each statement.
\placeholder{graph}{Red downward parabola with vertex (-1,4), x-intercepts -3 and 1, and y-intercept 3; arrows on both descending ends. Unit square grid, visible window approximately x=-5 to 3 and y=-3 to 5. Axes x/y arrowed; origin O; x labels -4, -2, 2 and y labels -2, 2, 4.}
\begin{tabular}{lll}
Statement & Yes & No \\
a. The domain is ((-\infty,4]). & blank & blank \\
b. The range is ((-\infty,4]). & blank & blank \\
c. The graph is positive on the interval ((0,\infty)). & blank & blank \\
d. The graph is decreasing on the interval ((-1,\infty)). & blank & blank \\
\end{tabular}
\answer{a. No; b. Yes; c. No; d. Yes.}
\end{practice}
\begin{practice}{32}{1}
SAT/ACT The graph shows the amount of money in an investment account. Which statement is true?
\placeholder{graph}{“Investment Account”: red V-shaped polyline from filled point (0,6) down to vertex (3,1), then up and right through approximately (7,8), with arrowhead on rising end. Horizontal x-axis “Year,” labels 0, 2, 4, 6; vertical y-axis “Amount (thousands of dollars),” labels 0, 2, 4, 6, 8. Unit grid; window x=0 to 7, y=0 to 9. Axes arrowed.}
A. \$6,000 was initially invested in the account.
B. \$1,000 was initially invested in the account.
C. At Year 3, there was \$0 in the account.
D. At Year 7, there was \$0 in the account.
\answer{A}
\end{practice}
\begin{practice}{33}{3}
Performance Task The graph shows the amount of water in a water tank over several hours.
\placeholder{graph}{“Amount of Water In Cistern”: red line segments join filled red points (0,12), (4,4), (6,4), (10,10). Horizontal x-axis “Hours” with labels 0, 2, 4, 6, 8, 10; vertical y-axis “Water (thousands of gallons)” with labels 0, 2, 4, 6, 8, 10, 12. Unit grid; window 0≤x≤10 and 0≤y≤12. Axes have arrowheads; polyline has no arrows.}
Part A What is the average rate of change on the interval ([0,4]) and on the interval ([6,10])? What is a possible explanation for what each rate of change indicates?
\answer{Part A (-2,000) gallons per hour; Over the interval ([0,4]), the water in the cistern is being used at a rate of 2,000 gallons per hour. 1,500 gallons per hour; Over the interval ([6,10]), the cistern is being filled with water at a rate of 1,500 gallons per hour.}
Part B What is a possible explanation for what occurred between 4 and 6 h?
\answer{Part B The water in the cistern is not being used, and the cistern is not being filled with water. Therefore, the amount of water in the cistern is constant.}
Part C What is the average rate of change on the interval ([0,10])? What does the rate of change mean? Does this rate of change give a good indication as to what is happening with the water in the cistern from 0 h to 10 h? Explain.
\answer{Part C (-200) gallons per hour; From 0 hours to 10 hours, the water in the cistern is being used at a rate of 200 gallons per hour. No, this rate of change shows the change in the starting amount of water at 0 hours and the ending amount of water at 10 hours, but it does not indicate that the amount of water decreased from 0 hours to 4 hours, remained the same from 4 hours to 6 hours, and increased from 6 hours to 10 hours.}
\end{practice}


% Source page 22: 12A STEP 4 Assess & Differentiate
\begin{te-addendum}{Assess}{RtISupport}
\textbf{LESSON QUIZ}
Use the Lesson Quiz to assess students’ understanding of the mathematics in the lesson.
Students can take the Lesson Quiz online or you can download a printable copy from SavvasRealize.com. The Lesson Quiz is also available in the Assessment Sourcebook.
\textbf{Item Analysis}
\begin{tabular}{lll}
Item & DOK & Skills Review and Practice \\
1 & 2 & F25 \\
2 & 2 & F25 \\
3 & 1 & F31 \\
4 & 2 & F25 \\
5 & 1 & F25 \\
\end{tabular}
Use the student scores on the Lesson Quiz to prescribe differentiated assignments.
If students take the Lesson Quiz online, it will be automatically scored and appropriate differentiated practice will be assigned based on student performance.
\begin{tabular}{lll}
Intervention & 0--3 points & Reteach to Build Understanding; Mathematical Literacy and Vocabulary; Lesson Virtual Nerd videos \\
On-Level & 4 points & Enrichment \\
Advanced & 5 points & Enrichment \\
\end{tabular}
\end{te-addendum}
\begin{te-addendum}{Assess}{GeneralizeETP}
\textbf{ASSESSMENT RESOURCES}
\textbf{1-1 Lesson Quiz: Key Features of Functions}
Name blank.
1. Benito produces and sells bulk coffee beans. The graph shows his profit as a function of the pounds of coffee beans sold. After selling a certain amount of coffee, Benito uses some earnings to buy more materials and then increases the price of the coffee. Which of the following is true?
\placeholder{graph}{Quiz 1 and 3: black polyline with filled dots at (0,-20), (1,-10), (2,0), (3,10), (4,20), (5,30), (6,40), (7,50), (8,25), (9,40), (10,55), (11,70). x-axis “Pounds of coffee sold,” labeled 2,4,6,8,10 with unit grid; y-axis “Profit (dollars),” labeled -20,20,40,60,80 with ten-dollar grid. Origin O, axes x/y with arrowheads, window approximately x=-2 to 12 and y=-40 to 90. No polyline arrows.}
A. Benito increases the price after selling 50 pounds of coffee beans.
B. The y-intercept represents the cost of one pound of coffee beans.
C. The x-intercept represents the number of pounds of coffee he has to sell before making a profit.
D. The domain of the function is (\{x\mid x\text{ is a real number greater than or equal to }0\}).
\answer{1. C}
2. The graph of the function (f(x)=-16x^2+10) shown at the right represents the distance, in feet, between a diver and th surface of the water, where (x) is the number of seconds after the diver leaves the diving board. What is the height of the diving board?
\te{Printed “th surface”; likely “the surface.”}
\placeholder{graph}{Quiz 2: black downward-curving arc f(x)=-16x^2+10 in first quadrant, begins at (0,10) and ends with downward-right arrow at approximately (0.79,0). x-axis “Time (s)” labels 0.5 and 1, grid 0.1; y-axis “Distance (ft)” labels 5 and 10, grid 1. Origin O, axes arrowed; visible window [0,1] by [0,10]. No point circles.}
\answer{2. 10 feet}
3. Which of the following correctly compares the two parts of the function in Item 1, over the intervals ([0,7]) and ([8,11])?
A. Benito made an extra \$25 after selling 7 pounds of coffee.
B. Benito spent \$1 after selling 50 pounds of coffee.
C. Benito’s profit per pound of coffee decreases from \$50 to \$25.
D. Benito’s profit per pound of coffee increases from \$10 to \$15.
\answer{3. D}
Use the graph shown.
\placeholder{graph}{Quiz 4–5: black downward parabola with vertex (1,4), x-intercepts -1 and 3, y-intercept 3, downward arrows at both ends. Unit square grid approximately x=-3 to 5, y=-6 to 6. Axes x/y arrowed, origin O; x labels -2,2,4, y labels -4,-2,2,4. No marked point circles.}
4. Identify the intervals on which the function is increasing and on which the function is decreasing.
Intervals on which the function increases: ((-\infty,\text{blank})).
Intervals on which the function decreases: ((\text{blank},\infty)).
\answer{4. Increases: ((-\infty,1)); decreases: ((1,\infty)).}
5. At what point does the function switch from increasing to decreasing? (blank, blank)
\answer{5. (1,4)}
\te{Lesson Quiz is available at SavvasRealize.com. enVision Algebra 2 • Assessment Sourcebook. Copyright © Savvas Learning Company LLC. All Rights Reserved.}
\end{te-addendum}

% Source page 23: 12B STEP 4 Assess & Differentiate
\begin{te-addendum}{Assess}{GeneralizeETP}
\textbf{DIFFERENTIATED RESOURCES}
I = Intervention; O = On-Level; A = Advanced.
\textbf{Reteach to Build Understanding—I}
Provides scaffolded reteaching for the key lesson concepts.
MathXL for School.
\textbf{1-1 Reteach to Build Understanding: Key Features of Functions}
Name blank.
Linear, quadratic, and absolute-value functions can be graphed using key features that the equation identifies. Key features can be labeled using either set builder notation or interval notation. Interval notation represents a set of real numbers by a pair of values. Parentheses are used when the points are not included. Brackets are used for when a boundary point is included.
1. Use the function (f(x)=x^2-16) and identify the key features that are listed in the table for the graph. Label the graph with the key features.
\begin{tabular}{ll}
Key Feature & Answer \\
Domain & \answer{((-\infty,\infty))} \\
Range & \answer{([-16,\infty))} \\
Minimum Value & \answer{(-16)} \\
x-intercept & \answer{(-4), 4} \\
y-intercept & \answer{(-16)} \\
Function is Positive & \answer{\uncertain{((-\infty,-4)) and ((4,\infty))}{Endpoint marks in the small preview may appear as brackets; inspect crop.}} \\
Function is Negative & \answer{((-4,4))} \\
Interval Increasing & \answer{((0,\infty))} \\
Interval Decreasing & \answer{((-\infty,0))} \\
\end{tabular}
\placeholder{graph}{Reteach item 1: black upward parabola f(x)=x^2-16 with vertex (0,-16), x-intercepts -4 and 4, arrows on upper ends. Grid approximately x=-5 to 5 in units of 1 and y=-25 to 25 in units of 5; x labels -4,-2,2,4; y labels -20,-10,10,20. Axes x/y arrowed, origin O.}
2. Cheyenne identified the key features of the function (f(x)=x+4) before graphing it. She wrote her answers in the table. Find and fix her error.
\begin{tabular}{ll}
Key Feature & Answer \\
Domain & ((-\infty,\infty)) \\
Range & ((-\infty,\infty)) \\
x intercept & (x=4) \answer{(x=-4)} \\
y intercept & 4 \\
Function is positive & ((-4,\infty)) \\
Function is negative & ((-\infty,-4)) \\
\end{tabular}
3. Identify the key features of the function (f(x)=|x-3|).
\begin{tabular}{lll}
Key Feature & Clues & Answer \\
Domain & Values of the x-axis & \answer{((-\infty,\infty))} \\
Range & Values of the y-axis & \answer{([0,\infty))} \\
Minimum value & Lowest point on the graph & 0 \\
x-intercept & When (y=0) & \answer{3} \\
y-intercept & When (x=0) & \answer{3} \\
Function is positive & Function is above the x axis & \answer{((-\infty,0)) and ((0,\infty))} \\
Function is negative & Function is below the x axis & NA \\
\end{tabular}
\te{Printed positive intervals exclude 0; likely they should exclude 3 for (f(x)=|x-3|).}
\end{te-addendum}
\begin{te-addendum}{Assess}{GeneralizeETP}
\textbf{Additional Practice—I, O}
Provides extra practice for each lesson.
MathXL for School.
\textbf{1-1 Additional Practice: Key Features of Functions}
Name blank.
For Items 1--2, identify the following information for the function of each graph:
1.
\placeholder{graph}{Additional Practice item 1: black upward parabola with vertex (0,-25), x-intercepts -5 and 5, both upper ends arrowed. Axes x/y, origin O; x labels -8,-4,4,8 and grid spacing 2; y labels -40,-20,20,40 and grid spacing 10. Visible window approximately x=-10 to 10 and y=-50 to 50.}
\answer{Domain: ((-\infty,\infty)); Range: ([-25,\infty)); x-intercepts: (-5), 5; y-intercepts: (-25); Interval positive: ((-\infty,-5)) and ((5,\infty)); Interval negative: ((-5,5)); Interval increasing: ((0,\infty)); Interval decreasing: ((-\infty,0)); Average rate of change over ([-5,0]): (-5).}
2.
\placeholder{graph}{Additional Practice item 2: black upward V, vertex (-2,3), y-intercept 7, steepness 2 on each side, arrows on both upper ends. Unit grid, axes x/y and origin O; x labels -6,-4,-2,2 and y labels 2,4,6,8. Window approximately x=-6 to 3, y=-1 to 9.}
\answer{Domain: ((-\infty,\infty)); Range: ([3,\infty)); x-intercepts: none; y-intercepts: 7; Interval positive: ((-\infty,\infty)); Interval negative: none; Interval increasing: ((-2,\infty)); Interval decreasing: ((-\infty,-2)); Average rate of change over ([-2,0]): 2.}
3. Sketch a linear graph given the following key features:
Domain: ((-\infty,\infty))
Range: ((-\infty,\infty))
Increasing: ((-\infty,\infty))
x-intercepts: (-2.5)
y-intercept: 10
Positive: ((-2.5,\infty))
Negative: ((-\infty,-2.5))
\answer{See graph.}
\placeholder{graph}{Additional Practice item 3 answer: magenta rising line with x-intercept -2.5 and y-intercept 10, arrows both ends. x/y axes with origin O, x grid spacing 1 and labels -6,-4,-2,2; y grid spacing 2 and labels -8,-4,4,8; visible window approximately x=-7 to 3, y=-10 to 10.}
4. Chiang is filling a (50\text{ ft}^3) container with water at a rate of (0.5\text{ ft}^3/\text{min}). Interpret the key features for this situation.
\answer{Domain: ([0,100]); Range: ([0,50]); Increasing: ([0,100]); Decreasing: N/A; x-intercepts: 0; y-intercepts: 0; Positive: ((0,100]); Negative: N/A.}
\end{te-addendum}
\begin{te-addendum}{Assess}{GeneralizeETP}
\textbf{Enrichment—O, A}
Presents engaging problems and activities that extend the lesson concepts.
MathXL for School.
\textbf{1-1 Enrichment: Key Features of Functions}
Name blank.
Complete each table for the given function. Use the corresponding letter shown in each table to answer the question.
\begin{tabular}{lll}
(f(x)=|x|-3) & Increasing & Decreasing \\
Positive & A (3,\infty) & B (-\infty,-3) \\
Negative & E \uncertain{[0,3)}{(0,3)} & F \uncertain{(-3,0]}{(-3,0)} \\
\end{tabular}
\begin{tabular}{lll}
(f(x)=2x^2-162) & Increasing & Decreasing \\
Positive & C (9,\infty) & D (-\infty,-9) \\
Negative & G \uncertain{[0,9)}{(0,9)} & H \uncertain{(-9,0]}{(-9,0)} \\
\end{tabular}
\begin{tabular}{lll}
(f(x)=6-|4x|) & Increasing & Decreasing \\
Positive & I \uncertain{(-1.5,0]}{(-1.5,0)} & J \uncertain{[0,1.5)}{(0,1.5)} \\
Negative & M (-\infty,-1.5) & N (1.5,\infty) \\
\end{tabular}
\begin{tabular}{lll}
(f(x)=x^2+4x) & Increasing & Decreasing \\
Positive & K (0,\infty) & L (-\infty,-4) \\
Negative & O \uncertain{[-2,0)}{(-2,0)} & P \uncertain{(-4,-2]}{(-4,-2)} \\
\end{tabular}
\begin{tabular}{lll}
(f(x)=|2x|-15) & Increasing & Decreasing \\
Positive & Q (7.5,\infty) & R (-\infty,-7.5) \\
Negative & U \uncertain{[0,7.5)}{(0,7.5)} & V \uncertain{(-7.5,0]}{(-7.5,0)} \\
\end{tabular}
\begin{tabular}{lll}
(f(x)=-x^2+28x) & Increasing & Decreasing \\
Positive & S (0,14) & T (14,28) \\
Negative & W (-\infty,0) & Y (28,\infty) \\
\end{tabular}
What did one function say to the other function?
\begin{tabular}{llll}
S & T & A & Y \\
(0,14) & (14,28) & (3,\infty) & (28,\infty) \\
\end{tabular}
\begin{tabular}{llllllll}
P & O & S & I & T & I & V & E \\
\uncertain{(-4,-2]}{(-4,-2)} & \uncertain{[-2,0)}{(-2,0)} & (0,14) & \uncertain{(-1.5,0]}{(-1.5,0)} & (14,28) & \uncertain{(-1.5,0]}{(-1.5,0)} & \uncertain{(-7.5,0]}{(-7.5,0)} & \uncertain{[0,3)}{(0,3)} \\
\end{tabular}
\answer{STAY POSITIVE}
\te{Endpoint delimiters in the small Enrichment preview remain difficult to distinguish after 300-dpi rendering; marked individually rather than silently inferred from the functions.}
\end{te-addendum}
\begin{te-addendum}{Assess}{ELLAddendum}
\textbf{Mathematical Literacy and Vocabulary—I, O}
Helps students develop and reinforce understanding of key terms and concepts.
\textbf{1-1 Mathematical Literacy and Vocabulary: Key Features of Functions}
Name blank.
\textbf{Concept List}
\begin{tabular}{lllll}
(x=-2) & (x=2) & (y=-2) & (y=2) & (-\infty,\infty) \\
(-\infty,-2) & (-2,\infty) & (0,\infty) & (\infty,0) & 2 \\
(-2) & A point on the graph where (x=0) & & A point on the graph where (y=0) & \\
When (f(x)<0) & Point where the function changes from increasing to decreasing. & & & \\
\end{tabular}
Complete the vocabulary chart, given the equation (f(x)=|x+2|). Use the information in the Concept List to fill in the chart.
\begin{tabular}{lll}
Word or Word Phrase & Definition & Picture or Example \\
Domain & Set of input values, or x-values & \answer{1. ((-\infty,\infty))} \\
Range & Set of output values or y-values & \answer{2. ((0,\infty))} \\
Maximum Value & Point where the function changes from increasing to decreasing & NA \\
Minimum value & Point where the function changes from decreasing to increasing & \answer{3. (-2)} \\
x-intercept & \answer{4. A point on the graph where (y=0)} & \answer{5. (-2)} \\
y-intercept & \answer{6. A point on the graph where (x=0)} & \answer{7. 2} \\
Function is positive & When (f(x)>0) & \answer{8. ((-\infty,\infty))} \\
Function is negative & \answer{9. When (f(x)<0)} & NA \\
Interval increasing & (f(x)) goes up on a portion of its domain when viewed from left to right. & \answer{10. ((-2,\infty))} \\
Interval decreasing & (f(x)) goes down on a portion of its domain when viewed from left to right. & \answer{11. ((-\infty,-2))} \\
\end{tabular}
\te{Printed range ((0,\infty)), minimum value (-2), and positive interval ((-\infty,\infty)); likely range ([0,\infty)), minimum value 0, and positivity excluding (x=-2) for the printed function (f(x)=|x+2|). The printed answers are retained.}
\end{te-addendum}
\begin{te-addendum}{Assess}{GeneralizeETP}
\textbf{Digital Differentiated Resources}
The Reteach to Build Understanding, Additional Practice, and Enrichment activities are available as digital assignments. These activities are automatically assigned when students complete the lesson quiz online and are automatically scored.
\textbf{Digital activity screenshot}
Solve the system of equations by graphing.
(y=3x+4)
(y=-2x+4)
Use the graphing tool to graph the system.
\placeholder{graph}{Digital exercise blank Cartesian grid with x/y axes and arrowheads, origin O, unit grid lines, numbered -10,-8,-6,-4,-2,2,4,6,8,10 on both axes. Window [-10,10] by [-10,10]; no functions or points plotted. A Question Help menu overlays the upper-right quadrant.}
Click to enlarge graph.
Click the graph, choose a tool in the palette and follow the instructions to create your graph.
\te{Interface labels: Exit; Question Help; Help Me Solve This; View an Example; Video; Textbook; Glossary; Math Tools; Print; 1 part remaining; Clear All; Check Answer; Review progress; Question 5 of 14; Go; Back; Next. The resource previews bear “enVision Algebra 2 • Teaching Resources,” SavvasRealize.com, and a copyright notice.}
\end{te-addendum}

\end{document}
